problem:  
Let \((M,\mathcal{D},g,\mu)\) be a compact, equiregular sub‑Riemannian manifold with Carnot–Carathéodory distance \(d_{SR}\).  
Define the **horizontal \(L^2\)-Wasserstein distance** \(W_{2,H}\) between two absolutely continuous probability measures \(\nu_0,\nu_1\in \mathcal{P}_2(M)\) by the Benamou–Brenier functional
\[
W_{2,H}^2(\nu_0,\nu_1)=\inf\!\left\{ \int_0^1\!\!\int_M |v_t|^2\,d\nu_t\,dt :
  (\nu_t,v_t)\ \text{satisfy }\partial_t\nu_t+\operatorname{div}_H(\nu_t v_t)=0,\ v_t\in\Gamma(\mathcal{D}) \right\},
\]
where \(\operatorname{div}_H\) is the horizontal divergence with respect to \(\mu\) and \(v_t\) takes values only in the horizontal distribution \(\mathcal{D}\).

For \(W_{2,H}\)–geodesics \((\nu_t)_{t\in[0,1]}\), define the usual Boltzmann entropy functional
\[
\mathrm{Ent}(\nu_t) = \int_M \rho_t \log\rho_t\, d\mu, \quad \nu_t = \rho_t \mu.
\]

**Conjecture (Horizontal‑entropy rigidity):**  
Assume there exists \(K\in\mathbb{R}\) with the following property:
> For every absolutely continuous pair \((\nu_0,\nu_1)\), the curve \((\nu_t)\) realizing \(W_{2,H}(\nu_0,\nu_1)\) satisfies  
> \[
> \mathrm{Ent}(\nu_t)\le (1-t)\mathrm{Ent}(\nu_0)+t\,\mathrm{Ent}(\nu_1)-\frac{K}{2}\,t(1-t)\,W_{2,H}^2(\nu_0,\nu_1),\quad \forall t\in[0,1].
> \]

Does this synthetic “horizontal \(K\)‑convexity of entropy” force the underlying infinitesimal geometry to be Riemannian?  
Equivalently:  

> If a compact, equiregular sub‑Riemannian space satisfies entropy \(K\)-convexity with respect to \(W_{2,H}\), must every metric tangent cone be Euclidean (so that \(\mathcal{D}=TM\))?

If not, characterize all possible non‑Riemannian tangent structures compatible with \(W_{2,H}\)-entropy convexity (if any exist).  

Why is it a "good" problem:  
This formulation gives a **precise rigidity conjecture**: it asks whether a natural generalization of displacement \(K\)-convexity—restricted to horizontal, sub‑Riemannian transports—already enforces infinitesimal Hilbertianity, thereby prohibiting genuine non‑Riemannian sub‑Riemannian geometries from satisfying any positive synthetic curvature bound.  

It is conceptually simple (“does horizontal entropy convexity already force Euclidean tangents?”) yet deeply nontrivial: the modified transport metric \(W_{2,H}\) involves a constrained continuity equation requiring horizontal flows, merging tools from optimal transport, sub‑elliptic analysis, and differential geometry of distributions.  
A proof or counterexample would delineate precisely **where the RCD/CD formalism breaks** and could inaugurate a rigorous synthetic theory for non‑Hilbertian geometries.  
Thus the problem balances clarity and depth—a crisp, testable conjecture with the potential to unify or sharply separate Riemannian and sub‑Riemannian curvature concepts.